
\documentclass[11pt]{article}
\usepackage[utf8]{inputenc}
\usepackage[spanish]{babel}
\usepackage{amsmath}
\usepackage{geometry}
\geometry{margin=1in}

\title{Documentación del Notebook de Series de Tiempo}
\date{}

\begin{document}
\maketitle

\section{Objetivo}
El notebook presenta una introducción práctica al análisis de series de tiempo utilizando Python.

\section{Temas}
\begin{itemize}
\item Procesos AR(1)
\item Comparación AR(1) vs AR(2)
\item Residuos y diagnóstico
\item AIC y BIC
\item Modelos ARIMA
\item Funciones ACF y PACF
\end{itemize}

\section{Fórmulas Relevantes}

Modelo AR(1):
\[
X_t = c + \phi X_{t-1} + \varepsilon_t
\]

AIC:
\[
AIC = 2k - 2\ln(L)
\]

BIC:
\[
BIC = k\ln(n) - 2\ln(L)
\]

Modelo ARIMA$(p,d,q)$.

\section{Conclusión}
La selección de modelos se apoya en AIC/BIC, mientras que ACF y PACF ayudan a identificar la estructura temporal de la serie.
\end{document}
